% ============================================================
% CHƯƠNG 3: THIẾT KẾ HỆ THỐNG  (kiến trúc logic — không đi sâu hiện thực mã)
% ============================================================
\clearpage
\section{Thiết kế hệ thống}
\label{sec:system_design}

Chương này mô tả \textbf{kiến trúc logic}: hệ thống gồm những \textbf{thành phần} nào,
mỗi thành phần \textbf{đảm nhiệm} gì và \textbf{tương tác} với nhau ra sao.
Hình~\ref{fig:system_overview} là sơ đồ trung tâm.
Việc lựa chọn công nghệ, cấu trúc thư mục mã nguồn, tham số thuật toán và đoạn lệnh cụ thể
được tập trung ở Chương~\ref{sec:implementation_experiment} (Hiện thực và thực nghiệm).

% ------------------------------------------------------------
\subsection{Kiến trúc tổng quan}

Hệ thống tổ chức theo mô hình \textbf{phân tầng}: dữ liệu thô được đưa qua lớp xử lý và tích hợp,
tập trung trong kho, rồi được hai lớp tiêu thụ độc lập là \textit{phân tích/trình bày} và \textit{khai phá dữ liệu}.

\begin{figure}[H]
  \centering
  \resizebox{0.96\textwidth}{!}{%
  \begin{tikzpicture}[
    node distance=6mm,
    font=\sffamily\footnotesize,
    >={Latex[length=2.2mm]},
    arr/.style={->, thick, draw=black!50},
    t0/.style={
      rectangle, rounded corners=5pt,
      draw=orange!80!black, fill=orange!8,
      line width=0.9pt, align=center,
      inner sep=10pt, text width=10cm
    },
    t1/.style={
      rectangle, rounded corners=5pt,
      draw=blue!50!black, fill=blue!5,
      line width=0.9pt, align=center,
      inner sep=10pt, text width=10cm
    },
    t2/.style={
      rectangle, rounded corners=5pt,
      draw=violet!60!black, fill=violet!5,
      line width=0.9pt, align=center,
      inner sep=10pt, text width=10cm
    },
    t3/.style={
      rectangle, rounded corners=5pt,
      draw=teal!55!black, fill=teal!5,
      line width=0.9pt, align=left,
      inner sep=10pt, text width=4.5cm,
      minimum height=2.6cm
    },
    lbl/.style={
      font=\scriptsize\itshape, text=black!50,
      align=right, text width=2.2cm
    },
  ]
  
  %% ── Tầng 0: Nguồn dữ liệu ──
  \node[t0] (src) {%
    \textbf{Nguồn dữ liệu thô}\\[3pt]
    {\scriptsize
    Tập dữ liệu chuyến bay Hoa Kỳ năm 2015 (Kaggle).
    Gồm thông tin chuyến bay, hãng hàng không và sân bay.
    Khoảng 5{,}8 triệu bản ghi, lưu dạng CSV chưa chuẩn hóa.}%
  };
  
  %% ── Tầng 1: ETL ──
  \node[t1, below=8mm of src] (etl) {%
    \textbf{Tầng ETL --- Xử lý \& Tích hợp}\\[3pt]
    {\scriptsize
    Trích xuất từ CSV theo từng lô (chunk) để tiết kiệm bộ nhớ.
    Làm sạch: loại bỏ giá trị thiếu, chuẩn hóa kiểu dữ liệu, mã hóa khóa ngoại.
    Nạp dữ liệu đã xử lý vào kho theo luồng một chiều:
    thô $\rightarrow$ sạch $\rightarrow$ kho.}%
  };
  \draw[arr] (src.south) -- (etl.north);
  
  %% ── Tầng 2: Data Warehouse ──
  \node[t2, below=8mm of etl] (dw) {%
    \textbf{Kho Dữ liệu (Data Warehouse)}\\[3pt]
    {\scriptsize
    Lưu trữ tích hợp, tổ chức theo \textbf{mô hình ngôi sao (Star Schema)}.
    Bảng sự kiện trung tâm đo lường chuyến bay, liên kết bốn bảng chiều:
    thời gian, hãng hàng không, sân bay và nguyên nhân delay.
    Là điểm truy cập chung cho cả phân tích và khai phá.}%
  };
  \draw[arr] (etl.south) -- (dw.north);
  
  %% ── Tầng 3: đặt thủ công, căn giữa so với dw ──
  %% dash bên trái, mine bên phải, khoảng cách tổng = 4.5+4.5+1cm = 10cm = text width dw
  \node[t3, below=10mm of dw, xshift=-2.55cm] (dash) {%
    \textbf{Phân tích \& Trình bày}\\[3pt]
    {\scriptsize
    Giao diện tương tác đa chiều: KPI tổng quan,
    phân tích theo tuyến bay, nguyên nhân delay,
    phân bố theo không gian và thời gian.
    Tích hợp mô hình dự báo delay cho chuyến bay mới.}%
  };
  
  \node[t3, right=10mm of dash] (mine) {%
    \textbf{Khai phá Dữ liệu}\\[3pt]
    {\scriptsize
    Truy vấn mẫu huấn luyện từ kho dữ liệu.
    Huấn luyện và so sánh năm thuật toán học máy.
    Đánh giá theo AUC-ROC, F1, Log Loss;
    xuất kết quả phục vụ báo cáo.}%
  };
  
  %% Mũi tên từ DW xuống hai khối
  \draw[arr] (dw.south) -- ++(0,-4mm) -| (dash.north);
  \draw[arr] (dw.south) -- ++(0,-4mm) -| (mine.north);
  
  %% Mũi tên ngang bổ sung
  \draw[dashed, black!30, thick, <->,
        shorten >=1.5mm, shorten <=1.5mm]
    (dash.east) --
    node[above, font=\tiny\itshape, text=black!45] {chia sẻ kết quả}
    (mine.west);
  
  %% ── Background ──
  \begin{scope}[on background layer]
    \node[
      fill=gray!3, rounded corners=8pt, inner sep=14pt,
      fit=(src)(etl)(dw)(dash)(mine),
      draw=gray!20, dashed, line width=0.6pt
    ] {};
  \end{scope}
  
  \end{tikzpicture}%
  }
  \caption{Sơ đồ kiến trúc tổng quan hệ thống}
  \label{fig:system_overview}
  \end{figure}

\subsubsection{Các thành phần và vai trò}

\begin{description}
  \item[Nguồn dữ liệu thô.]
  Tập bản ghi chuyến bay (khối lượng lớn) kèm các bảng tra cứu hãng và sân bay, ở dạng tệp bảng phẳng chưa chuẩn hóa.
  Đây là đầu vào duy nhất cho toàn bộ pipeline nghiệp vụ.

  \item[Tầng ETL (trích---biến đổi---nạp).]
  Đọc dữ liệu từ nguồn, làm sạch và ánh xạ sang khóa chiều;
  gắn mỗi chuyến bay với thời gian, hãng, sân bay và phân loại nguyên nhân trễ theo quy ước nghiệp vụ.
  Đầu ra là dữ liệu sẵn sàng nạp vào kho.

  \item[Kho dữ liệu (Data Warehouse).]
  Lưu trữ tích hợp, đóng vai trò \textbf{điểm truy cập chung} cho phân tích đa chiều và cho việc chuẩn bị dữ liệu huấn luyện.
  Tổ chức theo \textbf{Star Schema} (một bảng sự kiện trung tâm và các bảng chiều).

  \item[Phân tích và trình bày.]
  Cho phép người dùng xem chỉ số tổng hợp, drill theo thời gian/địa lý/hãng, phân tích nguyên nhân trễ
  và (ở mức thiết kế) tích hợp kết quả dự báo từ mô hình để hỗ trợ ra quyết định.

  \item[Khai phá dữ liệu.]
  Lấy dữ liệu (hoặc mẫu) phù hợp từ kho, xây dựng đặc trưng và nhãn phân loại theo nguyên tắc không rò rỉ thông tin tương lai;
  huấn luyện, so sánh và đánh giá mô hình; đầu ra phục vụ báo cáo và có thể được lớp trình bày sử dụng.
\end{description}

% ------------------------------------------------------------
\subsection{Luồng dữ liệu và tương tác giữa các thành phần}

\textbf{Luồng chính (một chiều):} nguồn thô $\rightarrow$ ETL $\rightarrow$ kho.
Hai lớp \textit{Phân tích} và \textit{Khai phá} chỉ \textbf{đọc} từ kho trong phạm vi thiết kế này,
không ghi ngược bản ghi nghiệp vụ vào kho.

\textbf{Tương tác giữa hai lớp tiêu thụ:} không bắt buộc cho vận hành kho;
ở mức ứng dụng, chúng có thể \textbf{chia sẻ kết quả} (ví dụ hiển thị xác suất trễ từ mô hình trên giao diện phân tích)
mà không làm thay đổi lược đồ dữ liệu trong kho.

\textbf{Tương tác với nguồn:} chỉ tầng ETL trao đổi trực tiếp với dữ liệu thô; các lớp phía trên chỉ thấy dữ liệu đã qua chuẩn hóa trong kho.

% ------------------------------------------------------------
\subsection{Mô hình dữ liệu logic trong kho (Star Schema)}

Trong kho, mỗi chuyến bay là một dòng trong \textbf{bảng sự kiện}, gắn với các \textbf{khóa chiều}
(hãng, sân bay xuất phát/đến, ngày, loại nguyên nhân trễ được gán theo quy tắc nghiệp vụ)
và chứa các \textbf{đo lường} như độ trễ, thời gian bay, quãng đường, cờ hủy/chuyển hướng, nhãn trễ theo ngưỡng thời gian.

\textbf{Lý do chọn Star Schema} so với chuẩn hóa sâu hơn (snowflake): truy vấn phân tích thường chỉ cần
nối fact với từng chiều một bước, phù hợp công cụ lưu trữ nhẹ và quy mô chiều nhỏ;
đánh đổi là dư thừa có kiểm soát trên các bảng chiều.

\begin{table}[H]
\centering
\caption{Cấu trúc logic bảng sự kiện \texttt{fact\_flights} (các nhóm thuộc tính chính)}
\vspace{4pt}
\small
\begin{tabular}{|l|l|}
\hline
\textbf{Nhóm} & \textbf{Nội dung} \\ \hline
Khóa \& chiều & Khóa surrogate; FK tới hãng, sân bay đi/đến, ngày, nguyên nhân trễ \\ \hline
Nhận dạng chuyến & Số hiệu chuyến, tail number (dạng degenerate dimension) \\ \hline
Thời gian lịch trình & Giờ cất/hạ cánh dự kiến và thực tế (theo quy ước lưu trữ) \\ \hline
Đo trễ \& thành phần & Trễ khởi hành/đến; các thành phần trễ (NAS, hãng, tàu bay, an ninh, thời tiết) \\ \hline
Cờ nghiệp vụ & Hủy chuyến, chuyển hướng, nhãn trễ theo ngưỡng (ví dụ $\geq 15$ phút) \\ \hline
\end{tabular}
\label{tab:fact_flights_logical}
\end{table}

\begin{table}[H]
\centering
\caption{Các bảng chiều (vai trò logic)}
\begin{tabular}{|l|l|}
\hline
\textbf{Bảng chiều} & \textbf{Vai trò} \\ \hline
\texttt{dim\_airline} & Mã và tên hãng \\ \hline
\texttt{dim\_airport} & Sân bay, địa lý (thành phố, bang, quốc gia, tọa độ) \\ \hline
\texttt{dim\_date} & Phân cấp thời gian: năm, quý, tháng, tuần, thứ, cuối tuần \\ \hline
\texttt{dim\_delay\_reason} & Danh mục nguyên nhân trễ (mã, tên, nhóm) \\ \hline
\end{tabular}
\label{tab:dimensions_logical}
\end{table}

\subsection*{Thiết kế phân tích và khai phá}

\textbf{Phân tích OLAP (khái niệm):} từ kho, các truy vấn tổng hợp theo nhiều chiều
(thời gian, hãng, sân bay, nguyên nhân, khung giờ) phục vụ báo cáo định hướng;
chi tiết nhóm câu hỏi và cách ánh xạ sang SQL nằm ở phần hiện thực.

\textbf{Mô hình học có giám sát (khái niệm):} nhãn dự báo gắn với trễ đến theo ngưỡng;
vector đặc trưng chỉ gồm thông tin có thể biết trước thời điểm dự báo (tránh rò rỉ tương lai);
so sánh nhiều học máy và xử lý mất cân bằng lớp bằng trọng số --- các lựa chọn thuật toán và tham số cụ thể
được trình bày ở Chương~\ref{sec:implementation_experiment}.
